% ============================================================
%  Section: Basis Functions and Convergence
%  For inclusion in ACTSC 971 American Options project report
%  Covers Longstaff-Schwartz (2001) and Tsitsiklis-Van Roy (2001)
% ============================================================

\section{Basis Functions and Convergence}

Both papers require approximating a continuation value function at each
backward time step using a finite set of basis functions. The choice of
these functions is a critical implementation decision that affects both
accuracy and numerical stability. The two papers treat this choice
differently --- one empirically and one theoretically --- and together
they provide a complete picture of when and why regression-based simulation
methods work.

% ------------------------------------------------------------
\subsection{Basis Functions in Longstaff--Schwartz (2001)}
% ------------------------------------------------------------

In the LSM framework, basis functions are used to approximate the
conditional expectation function $F(\omega; t_k)$ --- the expected
discounted payoff from continuing to hold the option at exercise date
$t_k$. The formal justification is that when this conditional expectation
belongs to the $L^2$ space of square-integrable functions, that space is a
Hilbert space with a countable orthonormal basis, so any such function can
be represented as a linear combination of basis elements \citep{longstaff2001}.

For the vanilla American put, the paper uses a constant plus the first
three weighted Laguerre polynomials:
\begin{align}
    L_0(X) &= e^{-X/2}, \label{eq:L0}\\
    L_1(X) &= e^{-X/2}(1 - X), \label{eq:L1}\\
    L_2(X) &= e^{-X/2}\!\left(1 - 2X + \tfrac{X^2}{2}\right), \label{eq:L2}
\end{align}
where $X = S_t / K$ is the stock price normalised by the strike. The
conditional expectation function is then approximated as
\begin{equation}
    F(\omega;\, t_{K-1}) \;=\; \sum_{j=0}^{\infty} a_j\, L_j(X),
    \label{eq:LSM_expansion}
\end{equation}
with coefficients $a_j$ estimated by ordinary least squares on the
in-the-money paths at each exercise date.

For the American--Bermuda--Asian (ABA) option with two state variables
$(S_t, A_t)$, the basis expands to eight functions: the first two Laguerre
polynomials in each of $S_t/K$ and $A_t/K$, together with their cross
products up to third order. For the five-asset maximum option in Section~8
of the paper, nineteen basis functions are used: a constant, the first five
Hermite polynomials in the maximum asset value, the values and squares of
the second through fifth ranked assets, all pairwise products of adjacent
ranked assets, and the product of all five assets.

A central empirical finding in \citet{longstaff2001} is that basis function
choice has surprisingly little effect on pricing accuracy. Section~8.3 of
the paper reports that replacing weighted Laguerre polynomials with simple
powers $S,\, S^2,\, S^3$, Hermite polynomials, or trigonometric functions
yields virtually identical results for the vanilla put. The explanation lies
in the nature of the regression: since LSM focuses on the \emph{fitted
value} rather than individual coefficients, multicollinearity among basis
functions does not bias prices, even if it makes individual coefficient
estimates unstable. The paper does caution that a nearly singular
cross-moment matrix can cause numerical problems, and recommends normalising
inputs --- for example dividing all prices by the strike $K$ --- to avoid
scaling issues.

% ------------------------------------------------------------
\subsection{Basis Functions in Tsitsiklis--Van Roy (2001)}
% ------------------------------------------------------------

In the TVR framework, basis functions serve a conceptually similar but
theoretically more precisely characterised role. Rather than approximating a
conditional expectation function directly, they parameterise the
$Q$-function (continuation value) as a linear combination:
\begin{equation}
    \tilde{Q}(x,\, r) \;=\; \sum_{k=1}^{K} r(k)\,\phi_k(x)
    \;=\; (\Phi r)(x),
    \label{eq:TVR_parameterisation}
\end{equation}
where $\phi_1, \ldots, \phi_K$ are fixed basis functions and $r \in
\mathbb{R}^K$ is the parameter vector estimated by regression at each
backward step \citep{tsitsiklis2001}.

The theoretical framework is deliberately agnostic about \emph{which}
basis functions are used. What matters is the quantity
\begin{equation}
    \epsilon_n^* \;=\; \min_{r}\,
    \bigl\|\Phi r - Q_n\bigr\|_{\pi_n}
    \;=\; \bigl\|\Pi_n Q_n - Q_n\bigr\|_{\pi_n},
    \label{eq:best_approx_error}
\end{equation}
the best approximation error achievable by the chosen basis at time step
$n$, measured in the norm induced by the trajectory distribution $\pi_n$.
As established in Theorem~1 of \citet{tsitsiklis2001}, the algorithm's
total approximation error is bounded in terms of $\epsilon_n^*$ (see
Section~\ref{sec:convergence} below), so any basis that makes $\epsilon_n^*$
small will produce a well-behaved algorithm. The LSM paper's empirical
finding that many different bases give similar results is consistent with
this: in one dimension, low-order polynomial and orthogonal polynomial
bases all achieve small $\epsilon^*$ for the smooth continuation value
functions that arise in standard option pricing problems.

The practical implication of TVR's framework becomes most important in
higher dimensions. In the ABA option, our LSM implementation uses the
eight-function Laguerre basis from the paper, while our TVR-style
comparator uses a simpler six-term polynomial basis
$[1,\, s,\, s^2,\, a,\, a^2,\, sa]$ where $s = S_t/K$ and $a = A_t/K$.
In TVR's framework, the polynomial basis has a larger $\epsilon^*$ than
the Laguerre basis --- it approximates the true continuation function less
accurately over the relevant state space. As shown in
Section~\ref{sec:results}, this larger $\epsilon^*$ feeds directly into
the error recursion, causing approximation errors to accumulate as the
backward induction proceeds, and producing the systematic overpricing
observed in the TVR-style comparator results.

% ------------------------------------------------------------
\subsection{Convergence}
\label{sec:convergence}
% ------------------------------------------------------------

The two papers offer fundamentally different types of convergence results,
reflecting their different goals. LSM provides practical convergence
guarantees suited to implementation, while TVR provides theoretical error
bounds that explain \emph{why} trajectory-based regression methods work.

\subsubsection{Convergence in Longstaff--Schwartz (2001)}

LSM convergence is established through two propositions. The first
addresses the bias of the algorithm and holds even when the option is
continuously exercisable.

\begin{proposition}[LSM downward bias; \citealp{longstaff2001}]
\label{prop:LSM_bias}
For any finite choice of $M$ basis functions, $K$ exercise dates, and
coefficient vector $\theta \in \mathbb{R}^{M \times (K-1)}$, the
following inequality holds almost surely:
\begin{equation}
    V(X) \;\geq\; \lim_{N \to \infty} \frac{1}{N}
    \sum_{i=1}^{N} \mathrm{LSM}(\omega_i;\, M, K).
    \label{eq:LSM_bias}
\end{equation}
\end{proposition}

The intuition is straightforward: the LSM algorithm produces a specific
stopping rule, and the true option value $V(X)$ is the supremum over
\emph{all} stopping rules. Any particular stopping rule --- including the
LSM one --- therefore gives a value no greater than the true optimum.
This result is important practically: LSM prices are conservative lower
bounds on the true American option value, never overstating it in the
limit.

The second proposition establishes full convergence for the single
exercise date case.

\begin{proposition}[LSM convergence; \citealp{longstaff2001}]
\label{prop:LSM_convergence}
Assume the option depends on a single state variable $X$ with support
$(0, \infty)$ following a Markov process, and can be exercised at two
dates $t_1 < t_2$. If the conditional expectation function $F(\omega; t_1)$
satisfies
\[
    \int_0^\infty e^{-X} F^2(\omega;\, t_1)\, dX < \infty
    \qquad \text{and} \qquad
    \int_0^\infty e^{-X} F_X^2(\omega;\, t_1)\, dX < \infty,
\]
then for any $\varepsilon > 0$ there exists $M < \infty$ such that
\begin{equation}
    \lim_{N \to \infty} \Pr\!\left(
        \left| V(X) - \frac{1}{N} \sum_{i=1}^{N}
        \mathrm{LSM}(\omega_i;\, M, K) \right| > \varepsilon
    \right) = 0.
    \label{eq:LSM_convergence}
\end{equation}
\end{proposition}

This result means that by choosing $M$ sufficiently large and letting
$N \to \infty$, the LSM algorithm converges to any desired degree of
accuracy. The key is uniform convergence of the fitted conditional
expectation $\hat{F}^M(\omega; t_1)$ to the true function $F^M(\omega;
t_1)$, which bounds the maximum pricing error. The paper conjectures that
analogous results hold in higher-dimensional settings.

\citet{longstaff2001} also validate convergence empirically through an
out-of-sample diagnostic test (Table~2 of their paper): regression
coefficients estimated on one set of paths are applied to a fresh set of
paths with different random seeds. The in-sample and out-of-sample values
are virtually identical across five seeds, with fewer than 5\% of
differences exceeding two standard errors, providing strong empirical
support for the stability of the stopping rule.

\subsubsection{Convergence in Tsitsiklis--Van Roy (2001)}

TVR's convergence analysis is more theoretically detailed, and
distinguishes between two separate questions: the quality of an idealised
algorithm with exact expectations and projections, and the convergence of
the practical sample-based algorithm to the idealised one.

\paragraph{Error blow-up with arbitrary representative states.}
A key negative result in \citet{tsitsiklis2001} is that naive approximate
value iteration can produce errors that grow \emph{exponentially} in the
problem horizon when representative states are chosen arbitrarily. In their
Section~IV example, an American put is approximated using a single basis
function $\phi(x) = g(x)$ and a projection measure $\pi$ concentrated at
a single fixed point $1 - \Delta$. The resulting parameter sequence
satisfies
\begin{equation}
    r_n \;\geq\; \frac{\alpha}{\Delta}(1-p)(1-d)\, r_{n+1},
    \label{eq:TVR_blowup}
\end{equation}
so that when $\Delta < \alpha(1-p)(1-d)$, the parameters grow at an
exponential rate as $n$ decreases from $N$ to $0$. Since the true option
value is bounded above by $\alpha^N$, this exponential growth implies that
the approximation error diverges. The root cause is that the representative
state $1-\Delta$ is far from where the risk-neutral process actually
visits, so errors compound at each backward step.

\paragraph{Bounded error under trajectory distributions.}
The remedy is to choose representative states by simulating forward
trajectories of the underlying process under the risk-neutral measure.
Let $\pi_n$ denote the true risk-neutral distribution of $x_n$. The
critical property that trajectory distributions possess is that the
transition operator $P$ is non-expansive across them. Specifically, by
Jensen's inequality and the tower property of conditional expectations:
\begin{equation}
    \|PJ\|_{\pi_n}^2
    \;=\; E\!\left[(PJ)(x_n)^2\right]
    \;\leq\; E\!\left[E\!\left[J^2(x_{n+1})\,\big|\,x_n\right]\right]
    \;=\; E\!\left[J^2(x_{n+1})\right]
    \;=\; \|J\|_{\pi_{n+1}}^2,
    \label{eq:TVR_nonexpansion}
\end{equation}
so $\|PJ\|_{\pi_n} \leq \|J\|_{\pi_{n+1}}$. This property ensures that
neither $P$ nor the projection operator $\Pi_n$ amplifies errors between
time steps. The main convergence result follows.

\begin{theorem}[TVR error bound; \citealp{tsitsiklis2001}]
\label{thm:TVR_theorem1}
Let $\pi_0, \ldots, \pi_{N-1}$ be the risk-neutral trajectory
distributions of $x_0, \ldots, x_{N-1}$. Define
$\tilde{e}_n = \|\tilde{Q}(\cdot, r_n) - Q_n\|_{\pi_n}$ as the
approximation error of the idealised algorithm, and
$\epsilon_n^* = \|\Pi_n Q_n - Q_n\|_{\pi_n}$ as the best achievable
error under the chosen basis. Then:
\begin{align}
    \tilde{e}_{N-1} &= \epsilon_{N-1}^*, \label{eq:TVR_init}\\
    \tilde{e}_n^2   &\leq \alpha^2\, \tilde{e}_{n+1}^2
                         + \left(\epsilon_n^*\right)^2,
    \qquad n = 0, 1, \ldots, N-2. \label{eq:TVR_recursion}
\end{align}
\end{theorem}

Unrolling the recursion in \eqref{eq:TVR_recursion}, if $\epsilon_n^* \leq
\epsilon$ uniformly across all time steps, the total error satisfies:
\begin{equation}
    \tilde{e}_n \;\leq\; \frac{\epsilon}{\sqrt{2(1 - \alpha^2)}},
    \label{eq:TVR_uniform_bound}
\end{equation}
which is \emph{independent of the time horizon $N$}. This is the central
theoretical achievement of \citet{tsitsiklis2001}: provided the basis is
good (small $\epsilon^*$) and states are sampled from the trajectory
distribution, the approximation error does not grow with the number of
time steps, in stark contrast to the exponential blow-up demonstrated in
Section~IV of their paper. Even the looser bound $\tilde{e}_n \leq
\sqrt{N}\,\epsilon$ grows only polynomially in $N$, not exponentially.

\paragraph{Convergence of the sample-based algorithm.}
Theorem~1 characterises the idealised algorithm. The practical algorithm
replaces exact expectations and projections with sample-based estimates
from simulated trajectories. Theorem~2 of \citet{tsitsiklis2001} addresses
this gap.

\begin{theorem}[Sample-based convergence; \citealp{tsitsiklis2001}]
\label{thm:TVR_theorem2}
As the number of simulated trajectories $m \to \infty$, the sample-based
parameter vectors $\hat{r}_n(m)$ converge almost surely to the parameter
vectors $r_n$ of the idealised algorithm:
\begin{equation}
    \hat{r}_n(m) \xrightarrow{\,a.s.\,} r_n
    \qquad \text{as } m \to \infty.
    \label{eq:TVR_as_convergence}
\end{equation}
\end{theorem}

The proof applies the strong law of large numbers to the least squares
objective: as $m$ grows, sample averages converge almost surely to their
population counterparts, so the sample regression converges to the
population regression. This provides the theoretical justification for
using more simulated paths to improve pricing accuracy.

% ------------------------------------------------------------
\subsection{Connecting the Two Papers}
% ------------------------------------------------------------

Reading the convergence results of both papers together yields a complete
and complementary picture. Proposition~\ref{prop:LSM_bias} from
\citet{longstaff2001} establishes that LSM prices are conservative ---
they never overstate the true American option value in the limit.
Theorem~\ref{thm:TVR_theorem1} from \citet{tsitsiklis2001} establishes
that when trajectory-based sampling is used, the approximation error
remains bounded regardless of horizon. Together, these results imply that
trajectory-based regression methods produce prices that are both bounded
above by the true option value and bounded below by a quantity that depends
only on the quality of the chosen basis functions, not on the number of
time steps.

The empirical results of the present study are consistent with both
theoretical frameworks. Our LSM prices sit close to but below the
binomial American benchmark across all strikes, consistent with the
downward bias property of Proposition~\ref{prop:LSM_bias}. Our TVR-style
comparator, by contrast, systematically overshoots the benchmark. This
overshoot is consistent with what Theorem~\ref{thm:TVR_theorem1} predicts
when $\epsilon^*$ is larger: the polynomial basis used in the TVR
comparator approximates the continuation function less accurately than the
Laguerre basis used in LSM (larger $\epsilon^*$), and the all-path
regression without an ITM filter produces a less focused approximation in
the region where the exercise decision matters most. The result is that
approximation errors accumulate more aggressively through the backward
recursion, producing the systematic overpricing observed across both the
vanilla put and the ABA option in our experiments. This empirical finding
directly illustrates the theoretical mechanism described in TVR
Section~IV: when the effective approximation error $\epsilon^*$ is large,
errors compound at each backward step rather than remaining bounded.
